\documentclass[12pt,a4paper]{article}
\usepackage{fontspec}\setmainfont{Latin Modern Roman}\setsansfont{DejaVu Sans}
\usepackage{amsmath,amssymb,graphicx,geometry,enumitem,longtable,booktabs,fancyhdr,hyperref}
\geometry{left=28mm,right=23mm,top=24mm,bottom=24mm}
\hypersetup{unicode=true,colorlinks=true,linkcolor=black,urlcolor=blue,pdftitle={Giải tích Chương 2 Mục 2.3 Tính giới hạn bằng các quy tắc giới hạn}}
\setlength{\parindent}{1.2em}\setlength{\parskip}{4pt}\setlength{\headheight}{16pt}
\pagestyle{fancy}\fancyhf{}\fancyhead[L]{\small Chương 2. Giới hạn và đạo hàm}\fancyhead[R]{\small §2.3}\fancyfoot[C]{\thepage}
\newcommand{\fig}[3][.68]{\begin{center}\begin{minipage}{\linewidth}\centering\includegraphics[width=#1\linewidth]{assets/#2.png}\\{\small #3}\end{minipage}\end{center}}
\newcommand{\ex}[1]{\par\medskip\noindent\textbf{Ví dụ #1.}\ }
\newcommand{\sol}{\par\noindent\textit{Lời giải.}\ }
\newcommand{\bt}[1]{\par\smallskip\noindent\textbf{#1.}\ }
\newcommand{\group}[1]{\par\medskip\noindent\textbf{#1}\par}
\newcommand{\gint}[1]{\left[\!\left[#1\right]\!\right]}
\begin{document}
\begin{center}
{\Large\bfseries GIẢI TÍCH}\\[5pt]
{\large Chương 2. Giới hạn và đạo hàm}\\[9pt]
{\Large\bfseries 2.3. Tính giới hạn bằng các quy tắc giới hạn}
\end{center}
\noindent\textbf{Nguồn:} James Stewart, Daniel K. Clegg, Saleem Watson, \textit{Calculus: Early Transcendentals}, ấn bản thứ 9 (2020), từ cuối trang in 94 đến hết trang 104 (trang PDF 129–139).
\section*{Các tính chất của giới hạn}
Ở §2.2, ta dùng máy tính và đồ thị để dự đoán giới hạn, nhưng đã thấy những phương pháp ấy không phải lúc nào cũng cho kết quả đúng. Trong mục này, ta dùng những tính chất sau, gọi là \textbf{các quy tắc giới hạn}, để tính giới hạn.
\subsection*{Các quy tắc giới hạn}
Giả sử $c$ là hằng số và hai giới hạn $\lim_{x\to a}f(x)$, $\lim_{x\to a}g(x)$ tồn tại. Khi đó:
\group{1. Quy tắc tổng}
\[\lim_{x\to a}[f(x)+g(x)]=\lim_{x\to a}f(x)+\lim_{x\to a}g(x).\]
\group{2. Quy tắc hiệu}
\[\lim_{x\to a}[f(x)-g(x)]=\lim_{x\to a}f(x)-\lim_{x\to a}g(x).\]
\group{3. Quy tắc nhân với hằng số}
\[\lim_{x\to a}[cf(x)]=c\lim_{x\to a}f(x).\]
\group{4. Quy tắc tích}
\[\lim_{x\to a}[f(x)g(x)]=\left(\lim_{x\to a}f(x)\right)\left(\lim_{x\to a}g(x)\right).\]
\group{5. Quy tắc thương}
\[\lim_{x\to a}\frac{f(x)}{g(x)}=\frac{\lim_{x\to a}f(x)}{\lim_{x\to a}g(x)},\qquad\text{nếu }\lim_{x\to a}g(x)\ne0.\]
Có thể phát biểu bằng lời như sau:
\begin{enumerate}
\item Giới hạn của một tổng bằng tổng các giới hạn.
\item Giới hạn của một hiệu bằng hiệu các giới hạn.
\item Giới hạn của tích một hằng số với hàm số bằng hằng số ấy nhân giới hạn của hàm số.
\item Giới hạn của một tích bằng tích các giới hạn.
\item Giới hạn của một thương bằng thương các giới hạn, với điều kiện giới hạn của mẫu khác 0.
\end{enumerate}
Dễ tin rằng các tính chất này đúng. Chẳng hạn, nếu $f(x)$ gần $L$ và $g(x)$ gần $M$, hợp lý khi kết luận $f(x)+g(x)$ gần $L+M$. Đó là cơ sở trực giác của Quy tắc 1. Ở §2.4, ta sẽ dùng định nghĩa chính xác để chứng minh quy tắc ấy. Các chứng minh còn lại được trình bày trong Phụ lục F.
\ex{1} Dùng các quy tắc giới hạn và đồ thị $f,g$ ở Hình 1 để tính, nếu tồn tại:
\[(a)\ \lim_{x\to-2}[f(x)+5g(x)];\qquad(b)\ \lim_{x\to1}[f(x)g(x)];\qquad(c)\ \lim_{x\to2}\frac{f(x)}{g(x)}.\]
\fig[.52]{f01}{Hình 1. Đồ thị các hàm $f$ và $g$ trong Ví dụ 1.}
\sol (a) Từ đồ thị, $\lim_{x\to-2}f(x)=1$, $\lim_{x\to-2}g(x)=-1$. Do đó,
\begin{align*}
\lim_{x\to-2}[f(x)+5g(x)]&=\lim_{x\to-2}f(x)+\lim_{x\to-2}[5g(x)]&&\text{(Quy tắc 1)}\\
&=\lim_{x\to-2}f(x)+5\lim_{x\to-2}g(x)&&\text{(Quy tắc 3)}\\
&=1+5(-1)=-4.
\end{align*}
(b) Ta thấy $\lim_{x\to1}f(x)=2$, nhưng $\lim_{x\to1}g(x)$ không tồn tại vì
\[\lim_{x\to1^-}g(x)=-2,\qquad\lim_{x\to1^+}g(x)=-1.\]
Không thể dùng Quy tắc 4 cho giới hạn hai phía đang xét, nhưng có thể dùng cho từng giới hạn một phía:
\[\lim_{x\to1^-}[f(x)g(x)]=\left(\lim_{x\to1^-}f(x)\right)\left(\lim_{x\to1^-}g(x)\right)=2(-2)=-4,\]
\[\lim_{x\to1^+}[f(x)g(x)]=\left(\lim_{x\to1^+}f(x)\right)\left(\lim_{x\to1^+}g(x)\right)=2(-1)=-2.\]
Hai giới hạn một phía khác nhau, nên giới hạn hai phía không tồn tại.

(c) Đồ thị cho $\lim_{x\to2}f(x)\approx1.4$ và $\lim_{x\to2}g(x)=0$. Không thể dùng Quy tắc 5. Giới hạn đã cho không tồn tại vì mẫu tiến đến 0, còn tử tiến đến một số khác 0.\hfill$\square$

Áp dụng nhiều lần quy tắc tích với $g(x)=f(x)$, ta được:
\group{6. Quy tắc lũy thừa}
\[\lim_{x\to a}[f(x)]^n=\left[\lim_{x\to a}f(x)\right]^n,\qquad n\text{ nguyên dương}.\]
Một tính chất tương tự đối với căn, được yêu cầu chứng minh trong Bài tập 2.5.69, là:
\group{7. Quy tắc căn}
\[\lim_{x\to a}\sqrt[n]{f(x)}=\sqrt[n]{\lim_{x\to a}f(x)},\qquad n\text{ nguyên dương}.\]
Nếu $n$ chẵn, giả sử $\lim_{x\to a}f(x)>0$.

Khi áp dụng bảy quy tắc này, ta cần hai giới hạn đặc biệt:
\group{8. Giới hạn của hằng số}
\[\lim_{x\to a}c=c.\]
\group{9. Giới hạn của hàm đồng nhất}
\[\lim_{x\to a}x=a.\]
Các giới hạn này hiển nhiên theo trực giác: hãy phát biểu bằng lời hoặc vẽ đồ thị $y=c$, $y=x$. Chứng minh bằng định nghĩa chính xác được yêu cầu trong Bài tập 2.4.23–24.

Đặt $f(x)=x$ trong Quy tắc 6 rồi dùng Quy tắc 9:
\group{10. Giới hạn của lũy thừa}
\[\lim_{x\to a}x^n=a^n,\qquad n\text{ nguyên dương}.\]
Đặt $f(x)=x$ trong Quy tắc 7 rồi dùng Quy tắc 9, ta được kết quả tương tự cho căn. Với căn bậc hai, chứng minh được phác thảo ở Bài tập 2.4.37.
\group{11. Giới hạn của căn}
\[\lim_{x\to a}\sqrt[n]{x}=\sqrt[n]{a},\qquad n\text{ nguyên dương}.\]
Nếu $n$ chẵn, giả sử $a>0$.
\ex{2} Tính các giới hạn sau và giải thích từng bước:
\[(a)\ \lim_{x\to5}(2x^2-3x+4);\qquad(b)\ \lim_{x\to-2}\frac{x^3+2x^2-1}{5-3x}.\]
\sol (a)
\begin{align*}
\lim_{x\to5}(2x^2-3x+4)&=\lim_{x\to5}(2x^2)-\lim_{x\to5}(3x)+\lim_{x\to5}4\\
&=2\lim_{x\to5}x^2-3\lim_{x\to5}x+\lim_{x\to5}4\\
&=2(5^2)-3(5)+4=39.
\end{align*}
Lần lượt dùng Quy tắc 2 và 1; Quy tắc 3; rồi Quy tắc 10, 9 và 8.

(b) Bắt đầu bằng Quy tắc 5; việc áp dụng chỉ được bảo đảm hoàn toàn ở bước cuối khi thấy tử và mẫu đều có giới hạn, và giới hạn mẫu khác 0:
\begin{align*}
\lim_{x\to-2}\frac{x^3+2x^2-1}{5-3x}
&=\frac{\lim_{x\to-2}(x^3+2x^2-1)}{\lim_{x\to-2}(5-3x)}\\
&=\frac{\lim_{x\to-2}x^3+2\lim_{x\to-2}x^2-\lim_{x\to-2}1}{\lim_{x\to-2}5-3\lim_{x\to-2}x}\\
&=\frac{(-2)^3+2(-2)^2-1}{5-3(-2)}=-\frac1{11}.
\end{align*}
Lần lượt dùng Quy tắc 5; Quy tắc 1, 2, 3; rồi Quy tắc 10, 9, 8.\hfill$\square$
\section*{Tính giới hạn bằng phép thế trực tiếp}
Trong Ví dụ 2(a), với $f(x)=2x^2-3x+4$, ta đã có $\lim_{x\to5}f(x)=39=f(5)$: chỉ cần thay 5 vào $x$ cũng được kết quả đúng. Tương tự ở (b). Hai hàm này lần lượt là đa thức và hàm hữu tỉ; áp dụng các quy tắc giới hạn theo cách tương tự chứng minh phép thế trực tiếp luôn đúng cho chúng (xem Bài tập 59, 60).
\begin{quote}\textbf{Tính chất thế trực tiếp.} Nếu $f$ là đa thức hoặc hàm hữu tỉ, và $a$ thuộc tập xác định của $f$, thì
\[\lim_{x\to a}f(x)=f(a).\]
\end{quote}
Hàm có tính chất thế trực tiếp được gọi là \textbf{liên tục tại $a$}, sẽ được nghiên cứu ở §2.5. Tuy nhiên, không phải giới hạn nào cũng tính được ngay bằng phép thế trực tiếp.
\subsection*{Newton và giới hạn}
Isaac Newton sinh vào ngày Giáng sinh năm 1642, năm Galileo qua đời. Khi vào Đại học Cambridge năm 1661, Newton chưa biết nhiều toán học; ông học nhanh qua các tác phẩm Euclid, Descartes và các bài giảng của Isaac Barrow. Cambridge đóng cửa vì dịch hạch trong các năm 1665–1666, nên Newton về quê suy ngẫm những điều đã học. Hai năm ấy đặc biệt hiệu quả: ông thực hiện bốn khám phá lớn, gồm biểu diễn hàm bằng tổng các chuỗi vô hạn, kể cả định lý nhị thức; nghiên cứu phép tính vi phân và tích phân; các định luật chuyển động và định luật vạn vật hấp dẫn; thí nghiệm lăng kính về bản chất ánh sáng, màu sắc.

Lo ngại tranh luận và chỉ trích, ông chần chừ công bố. Mãi năm 1687, theo sự thúc giục của nhà thiên văn Halley, Newton mới xuất bản \textit{Principia Mathematica}. Trong tác phẩm mà sách đánh giá là luận thuyết khoa học vĩ đại nhất từng được viết, Newton trình bày cách tiếp cận giải tích của mình, vận dụng để nghiên cứu cơ học, động lực học chất lưu, chuyển động sóng và giải thích chuyển động hành tinh, sao chổi.

Nguồn gốc giải tích có trong việc tính diện tích, thể tích của các học giả Hy Lạp cổ như Eudoxus và Archimedes. Ý tưởng giới hạn hàm chứa trong ``phương pháp vét kiệt'' của họ, nhưng chưa được phát biểu tường minh. Các tiền bối trực tiếp của Newton như Cavalieri, Fermat, Barrow cũng chưa thực sự dùng giới hạn. Newton là người đầu tiên nói tường minh về giới hạn, giải thích ý tưởng chính là các đại lượng ``tiến lại gần nhau hơn mọi độ chênh lệch cho trước''. Ông coi giới hạn là khái niệm cơ sở của giải tích; các nhà toán học về sau như Cauchy tiếp tục làm sáng tỏ ý tưởng này.
\ex{3} Tìm $\displaystyle\lim_{x\to1}\frac{x^2-1}{x-1}$.
\sol Đặt $f(x)=(x^2-1)/(x-1)$. Không thể thay $x=1$ vì $f(1)$ không xác định; cũng không thể áp dụng quy tắc thương vì giới hạn mẫu bằng 0. Ta cần biến đổi đại số trước, bằng cách phân tích tử thành hiệu hai bình phương:
\[\frac{x^2-1}{x-1}=\frac{(x-1)(x+1)}{x-1}.\]
Khi $x\to1$, ta xét $x\ne1$, nên $x-1\ne0$; có thể khử nhân tử chung rồi thế trực tiếp:
\[\lim_{x\to1}\frac{x^2-1}{x-1}=\lim_{x\to1}\frac{(x-1)(x+1)}{x-1}=\lim_{x\to1}(x+1)=1+1=2.\]
Giới hạn này đã xuất hiện ở Ví dụ 2.1.1 khi tìm tiếp tuyến parabol $y=x^2$ tại $(1,1)$.\hfill$\square$
\noindent\textit{Chú thích bên lề:} Dù mẫu tiến đến 0, Ví dụ 3 không có giới hạn vô cực. Khi cả tử và mẫu cùng tiến đến 0, giới hạn có thể vô cực hoặc hữu hạn.

\noindent\textbf{Lưu ý.} Ta đã thay $f(x)=(x^2-1)/(x-1)$ bằng hàm đơn giản hơn $g(x)=x+1$ có cùng giới hạn. Điều đó hợp lệ vì $f(x)=g(x)$ khi $x\ne1$, và giới hạn tại 1 không xét giá trị tại chính 1. Tổng quát:
\begin{quote}Nếu $f(x)=g(x)$ khi $x\ne a$, thì $\lim_{x\to a}f(x)=\lim_{x\to a}g(x)$, với điều kiện các giới hạn tồn tại.\end{quote}
\ex{4} Tìm $\lim_{x\to1}g(x)$ với
\[g(x)=\begin{cases}x+1,&x\ne1,\\\pi,&x=1.\end{cases}\]
\sol $g$ xác định tại 1 và $g(1)=\pi$, nhưng giới hạn không phụ thuộc giá trị tại 1. Vì $g(x)=x+1$ khi $x\ne1$,
\[\lim_{x\to1}g(x)=\lim_{x\to1}(x+1)=2.\qquad\square\]
Giá trị của hai hàm trong Ví dụ 3, 4 trùng nhau trừ tại $x=1$ (Hình 2); do đó giới hạn tại 1 bằng nhau.
\fig[.43]{f02}{Hình 2. Đồ thị $f$ ở Ví dụ 3 (trên) và $g$ ở Ví dụ 4 (dưới).}
\ex{5} Tính $\displaystyle\lim_{h\to0}\frac{(3+h)^2-9}{h}$.
\sol Với $F(h)=((3+h)^2-9)/h$, không thể cho $h=0$ vì $F(0)$ không xác định. Biến đổi đại số cho
\[F(h)=\frac{(9+6h+h^2)-9}{h}=\frac{6h+h^2}{h}=\frac{h(6+h)}h=6+h.\]
Nhớ rằng chỉ xét $h\ne0$. Do đó,
\[\lim_{h\to0}\frac{(3+h)^2-9}{h}=\lim_{h\to0}(6+h)=6.\qquad\square\]
\ex{6} Tìm $\displaystyle\lim_{t\to0}\frac{\sqrt{t^2+9}-3}{t^2}$.
\sol Không thể áp dụng ngay quy tắc thương vì giới hạn mẫu bằng 0. Ta biến đổi bằng cách hữu tỉ hóa tử:
\begin{align*}
\lim_{t\to0}\frac{\sqrt{t^2+9}-3}{t^2}
&=\lim_{t\to0}\frac{\sqrt{t^2+9}-3}{t^2}\cdot\frac{\sqrt{t^2+9}+3}{\sqrt{t^2+9}+3}\\
&=\lim_{t\to0}\frac{(t^2+9)-9}{t^2(\sqrt{t^2+9}+3)}\\
&=\lim_{t\to0}\frac{t^2}{t^2(\sqrt{t^2+9}+3)}\\
&=\lim_{t\to0}\frac1{\sqrt{t^2+9}+3}\\
&=\frac1{\left(\lim_{t\to0}\sqrt{t^2+9}\right)+3}
=\frac1{3+3}=\frac16.
\end{align*}
Ở bước cuối, dùng các tính chất 5, 1, 7, 8, 10. Phép tính xác nhận dự đoán trong Ví dụ 2.2.1.\hfill$\square$
\section*{Dùng giới hạn một phía}
Một số giới hạn nên tính bằng cách tìm giới hạn bên trái và bên phải trước. Định lý sau nhắc lại kết quả §2.2: giới hạn hai phía tồn tại khi và chỉ khi hai giới hạn một phía cùng tồn tại và bằng nhau.
\begin{quote}\textbf{1. Định lý.}
\[\lim_{x\to a}f(x)=L\quad\Longleftrightarrow\quad\lim_{x\to a^-}f(x)=L=\lim_{x\to a^+}f(x).\]
\end{quote}
Các quy tắc giới hạn cũng đúng đối với giới hạn một phía.
\ex{7} Chứng minh $\lim_{x\to0}|x|=0$.
\sol Nhắc lại $|x|=x$ khi $x\ge0$ và $|x|=-x$ khi $x<0$. Vì vậy,
\[\lim_{x\to0^+}|x|=\lim_{x\to0^+}x=0;\qquad\lim_{x\to0^-}|x|=\lim_{x\to0^-}(-x)=0.\]
Theo Định lý 1, $\lim_{x\to0}|x|=0$. Hình 3 cho thấy kết quả này hợp lý.
\fig[.45]{f03}{Hình 3. Đồ thị $y=|x|$. \(\square\)}
\ex{8} Chứng minh $\displaystyle\lim_{x\to0}\frac{|x|}x$ không tồn tại.
\sol Với $x>0$, $|x|=x$; với $x<0$, $|x|=-x$. Do đó,
\[\lim_{x\to0^+}\frac{|x|}x=\lim_{x\to0^+}\frac xx=\lim_{x\to0^+}1=1,\]
\[\lim_{x\to0^-}\frac{|x|}x=\lim_{x\to0^-}\frac{-x}x=\lim_{x\to0^-}(-1)=-1.\]
Hai giới hạn một phía khác nhau, nên giới hạn không tồn tại theo Định lý 1. Hình 4 phù hợp với các giới hạn vừa tìm được.
\fig[.45]{f04}{Hình 4. Đồ thị $y=|x|/x$. \(\square\)}
\ex{9} Cho $f(x)=\begin{cases}\sqrt{x-4},&x>4,\\8-2x,&x<4.\end{cases}$ Xác định xem $\lim_{x\to4}f(x)$ có tồn tại không.
\sol Vì $f(x)=\sqrt{x-4}$ khi $x>4$,
\[\lim_{x\to4^+}f(x)=\lim_{x\to4^+}\sqrt{x-4}=\sqrt{4-4}=0.\]
Vì $f(x)=8-2x$ khi $x<4$,
\[\lim_{x\to4^-}f(x)=\lim_{x\to4^-}(8-2x)=8-2\cdot4=0.\]
Hai giới hạn bằng nhau, nên giới hạn tồn tại và $\lim_{x\to4}f(x)=0$. Đồ thị ở Hình 5.
\noindent\textit{Chú thích của sách:} Ví dụ 2.4.4 chứng minh $\lim_{x\to0^+}\sqrt{x}=0$.
\fig[.48]{f05}{Hình 5. Đồ thị hàm ở Ví dụ 9. \(\square\)}
\ex{10} Hàm phần nguyên được định nghĩa bởi $\gint{x}$ bằng số nguyên lớn nhất không vượt quá $x$. Chẳng hạn,
\[\gint{4}=4,\quad\gint{4.8}=4,\quad\gint{\pi}=3,\quad\gint{\sqrt2}=1,\quad\gint{-\tfrac12}=-1.\]
Chứng minh $\lim_{x\to3}\gint{x}$ không tồn tại.
\sol Đồ thị ở Hình 6. Vì $\gint{x}=3$ khi $3\le x<4$,
\[\lim_{x\to3^+}\gint{x}=\lim_{x\to3^+}3=3.\]
Vì $\gint{x}=2$ khi $2\le x<3$,
\[\lim_{x\to3^-}\gint{x}=\lim_{x\to3^-}2=2.\]
Hai giới hạn khác nhau, nên giới hạn hai phía không tồn tại theo Định lý 1.
\fig[.49]{f06}{Hình 6. Hàm phần nguyên. \(\square\)}
\noindent\textit{Chú thích bên lề:} Các ký hiệu khác là $[x]$ và $\lfloor x\rfloor$. Hàm phần nguyên cũng được gọi là \textit{floor function}.
\section*{Định lý kẹp}
Hai định lý sau mô tả quan hệ giữa các giới hạn khi giá trị một hàm lớn hơn hoặc bằng giá trị hàm khác. Chứng minh có trong Phụ lục F.
\begin{quote}\textbf{2. Định lý.} Nếu $f(x)\le g(x)$ khi $x$ gần $a$, có thể trừ $a$, và cả hai hàm đều có giới hạn khi $x\to a$, thì
\[\lim_{x\to a}f(x)\le\lim_{x\to a}g(x).\]
\end{quote}
\begin{quote}\textbf{3. Định lý kẹp.} Nếu $f(x)\le g(x)\le h(x)$ khi $x$ gần $a$, có thể trừ $a$, và
\[\lim_{x\to a}f(x)=\lim_{x\to a}h(x)=L,\]
thì $\lim_{x\to a}g(x)=L$.
\end{quote}
Định lý kẹp, còn gọi là \textit{Sandwich Theorem} hoặc \textit{Pinching Theorem}, được minh họa ở Hình 7. Nếu $g(x)$ bị kẹp giữa $f(x),h(x)$ gần $a$, và hai hàm ngoài có cùng giới hạn $L$, thì $g$ buộc phải có giới hạn $L$ tại $a$.
\fig[.52]{f07}{Hình 7. Minh họa định lý kẹp.}
\ex{11} Chứng minh $\displaystyle\lim_{x\to0}x^2\sin\frac1x=0$.
\sol Không thể viết giới hạn thành tích của $\lim_{x\to0}x^2$ và $\lim_{x\to0}\sin(1/x)$, vì giới hạn thứ hai không tồn tại (xem Ví dụ 2.2.5).

Ta dùng định lý kẹp. Cần tìm $f$ nhỏ hơn $g(x)=x^2\sin(1/x)$ và $h$ lớn hơn $g$, sao cho cả hai tiến đến 0. Vì sin của mọi số nằm giữa $-1$ và 1,
\begin{equation}-1\le\sin\frac1x\le1.\tag{4}\end{equation}
Nhân một bất đẳng thức với số dương không đổi chiều bất đẳng thức. Vì $x^2\ge0$, nhân các vế của (4) với $x^2$ cho
\[-x^2\le x^2\sin\frac1x\le x^2.\]
Hình 8 minh họa điều này. Ta biết
\[\lim_{x\to0}x^2=0,\qquad\lim_{x\to0}(-x^2)=0.\]
Đặt $f(x)=-x^2$, $g(x)=x^2\sin(1/x)$, $h(x)=x^2$ trong định lý kẹp, ta được
\[\lim_{x\to0}x^2\sin\frac1x=0.\qquad\square\]
\fig[.52]{f08}{Hình 8. $y=x^2\sin(1/x)$ nằm giữa $y=-x^2$ và $y=x^2$.}
\section*{2.3. Bài tập}
\bt{1} Cho $\lim_{x\to2}f(x)=4$, $\lim_{x\to2}g(x)=-2$, $\lim_{x\to2}h(x)=0$. Tính những giới hạn tồn tại; nếu không tồn tại, giải thích:
\[(a)\ \lim_{x\to2}[f(x)+5g(x)];\quad(b)\ \lim_{x\to2}[g(x)]^3;\quad(c)\ \lim_{x\to2}\sqrt{f(x)};\]
\[(d)\ \lim_{x\to2}\frac{3f(x)}{g(x)};\qquad(e)\ \lim_{x\to2}\frac{g(x)}{h(x)};\qquad(f)\ \lim_{x\to2}\frac{g(x)h(x)}{f(x)}.\]
\bt{2} Cho đồ thị $f,g$. Dùng chúng tính mỗi giới hạn, nếu tồn tại; nếu không, giải thích:
\[(a)\ \lim_{x\to2}[f(x)+g(x)];\quad(b)\ \lim_{x\to0}[f(x)-g(x)];\quad(c)\ \lim_{x\to-1}[f(x)g(x)];\]
\[(d)\ \lim_{x\to3}\frac{f(x)}{g(x)};\quad(e)\ \lim_{x\to2}[x^2f(x)];\quad(f)\ f(-1)+\lim_{x\to-1}g(x).\]
\fig[.86]{e02}{Đồ thị cho Bài tập 2.}
\group{3–9. Tính giới hạn, giải thích từng bước bằng cách nêu quy tắc giới hạn thích hợp.}
\bt{3} $\displaystyle\lim_{x\to5}(4x^2-5x)$.
\bt{4} $\displaystyle\lim_{x\to-3}(2x^3+6x^2-9)$.
\bt{5} $\displaystyle\lim_{v\to2}(v^2+2v)(2v^3-5)$.
\bt{6} $\displaystyle\lim_{t\to7}\frac{3t^2+1}{t^2-5t+2}$.
\bt{7} $\displaystyle\lim_{u\to-2}\sqrt{9-u^3+2u^2}$.
\bt{8} $\displaystyle\lim_{x\to3}\sqrt[3]{x+5}\,(2x^2-3x)$.
\bt{9} $\displaystyle\lim_{t\to-1}\left(\frac{2t^5-t^4}{5t^2+4}\right)^3$.
\bt{10} (a) Đẳng thức sau sai ở điểm nào?
\[\frac{x^2+x-6}{x-2}=x+3.\]
(b) Dựa vào (a), giải thích vì sao đẳng thức sau đúng:
\[\lim_{x\to2}\frac{x^2+x-6}{x-2}=\lim_{x\to2}(x+3).\]
\group{11–34. Tính giới hạn, nếu tồn tại.}
\bt{11} $\displaystyle\lim_{x\to-2}(3x-7)$.
\bt{12} $\displaystyle\lim_{x\to6}(8-\tfrac12 x)$.
\bt{13} $\displaystyle\lim_{t\to4}\frac{t^2-2t-8}{t-4}$.
\bt{14} $\displaystyle\lim_{x\to-3}\frac{x^2+3x}{x^2-x-12}$.
\bt{15} $\displaystyle\lim_{x\to2}\frac{x^2+5x+4}{x-2}$.
\bt{16} $\displaystyle\lim_{x\to4}\frac{x^2+3x}{x^2-x-12}$.
\bt{17} $\displaystyle\lim_{x\to-2}\frac{x^2-x-6}{3x^2+5x-2}$.
\bt{18} $\displaystyle\lim_{x\to-5}\frac{2x^2+9x-5}{x^2-25}$.
\bt{19} $\displaystyle\lim_{t\to3}\frac{t^3-27}{t^2-9}$.
\bt{20} $\displaystyle\lim_{u\to-1}\frac{u+1}{u^3+1}$.
\bt{21} $\displaystyle\lim_{h\to0}\frac{(h-3)^2-9}{h}$.
\bt{22} $\displaystyle\lim_{x\to9}\frac{9-x}{3-\sqrt{x}}$.
\bt{23} $\displaystyle\lim_{h\to0}\frac{\sqrt{9+h}-3}{h}$.
\bt{24} $\displaystyle\lim_{x\to2}\frac{2-x}{\sqrt{x+2}-2}$.
\bt{25} $\displaystyle\lim_{x\to3}\frac{1/x-1/3}{x-3}$.
\bt{26} $\displaystyle\lim_{h\to0}\frac{(-2+h)^{-1}+2^{-1}}h$.
\bt{27} $\displaystyle\lim_{t\to0}\frac{\sqrt{1+t}-\sqrt{1-t}}t$.
\bt{28} $\displaystyle\lim_{t\to0}\left(\frac1t-\frac1{t^2+t}\right)$.
\bt{29} $\displaystyle\lim_{x\to16}\frac{4-\sqrt{x}}{16x-x^2}$.
\bt{30} $\displaystyle\lim_{x\to2}\frac{x^2-4x+4}{x^4-3x^2-4}$.
\bt{31} $\displaystyle\lim_{t\to0}\left(\frac1{t\sqrt{1+t}}-\frac1t\right)$.
\bt{32} $\displaystyle\lim_{x\to-4}\frac{\sqrt{x^2+9}-5}{x+4}$.
\bt{33} $\displaystyle\lim_{h\to0}\frac{(x+h)^3-x^3}{h}$.
\bt{34} $\displaystyle\lim_{h\to0}\frac{1/(x+h)^2-1/x^2}{h}$.
\bt{35} (a) Vẽ $f(x)=x/(\sqrt{1+3x}-1)$ để ước lượng
\[\lim_{x\to0}\frac{x}{\sqrt{1+3x}-1}.\]
(b) Lập bảng $f(x)$ khi $x$ gần 0 rồi dự đoán giới hạn. (c) Dùng các quy tắc giới hạn chứng minh dự đoán.
\bt{36} (a) Dùng đồ thị $f(x)=(\sqrt{3+x}-\sqrt3)/x$ ước lượng $\lim_{x\to0}f(x)$ đến hai chữ số sau dấu phẩy. (b) Dùng bảng giá trị ước lượng đến bốn chữ số sau dấu phẩy. (c) Dùng các quy tắc giới hạn tìm giá trị chính xác.
\bt{37} Dùng định lý kẹp chứng minh $\lim_{x\to0}x^2\cos20\pi x=0$. Minh họa bằng cách vẽ $f(x)=-x^2$, $g(x)=x^2\cos20\pi x$, $h(x)=x^2$ trên cùng màn hình.
\bt{38} Dùng định lý kẹp chứng minh $\displaystyle\lim_{x\to0^+}\sqrt{x^3+x^2}\sin\frac\pi x=0$. Minh họa bằng đồ thị $f,g,h$ (theo ký hiệu định lý kẹp) trên cùng màn hình.
\bt{39} Nếu $4x-9\le f(x)\le x^2-4x+7$ với $x\ge0$, tìm $\lim_{x\to4}f(x)$.
\bt{40} Nếu $2x\le g(x)\le x^4-x^2+2$ với mọi $x$, tính $\lim_{x\to1}g(x)$.
\bt{41} Chứng minh $\displaystyle\lim_{x\to0}x^4\cos\frac2x=0$.
\bt{42} Chứng minh $\displaystyle\lim_{x\to0^+}\sqrt{x}\,e^{\sin(\pi/x)}=0$.
\group{43–48. Tìm giới hạn, nếu tồn tại; nếu không, giải thích vì sao.}
\bt{43} $\displaystyle\lim_{x\to-4}(|x+4|-2x)$.
\bt{44} $\displaystyle\lim_{x\to-4}\frac{|x+4|}{2x+8}$.
\bt{45} $\displaystyle\lim_{x\to0.5^-}\frac{2x-1}{|2x^3-x^2|}$.
\bt{46} $\displaystyle\lim_{x\to-2}\frac{2-|x|}{2+x}$.
\bt{47} $\displaystyle\lim_{x\to0^-}\left(\frac1x-\frac1{|x|}\right)$.
\bt{48} $\displaystyle\lim_{x\to0^+}\left(\frac1x-\frac1{|x|}\right)$.
\bt{49} \textbf{Hàm dấu.} Hàm dấu, ký hiệu $\operatorname{sgn}$, được định nghĩa bởi
\[\operatorname{sgn}x=\begin{cases}-1,&x<0,\\0,&x=0,\\1,&x>0.\end{cases}\]
(a) Vẽ phác đồ thị. (b) Tìm các giới hạn hoặc giải thích vì sao không tồn tại:
\[(i)\ \lim_{x\to0^+}\operatorname{sgn}x;\quad(ii)\ \lim_{x\to0^-}\operatorname{sgn}x;\]
\[(iii)\ \lim_{x\to0}\operatorname{sgn}x;\qquad(iv)\ \lim_{x\to0}|\operatorname{sgn}x|.\]
\bt{50} Cho $g(x)=\operatorname{sgn}(\sin x)$.

(a) Tìm các giới hạn hoặc giải thích vì sao không tồn tại:
\[(i)\ \lim_{x\to0^+}g(x);\quad(ii)\ \lim_{x\to0^-}g(x);\quad(iii)\ \lim_{x\to0}g(x);\]
\[(iv)\ \lim_{x\to\pi^+}g(x);\quad(v)\ \lim_{x\to\pi^-}g(x);\quad(vi)\ \lim_{x\to\pi}g(x).\]
(b) Với những $a$ nào, $\lim_{x\to a}g(x)$ không tồn tại? (c) Vẽ phác đồ thị $g$.
\bt{51} Cho $g(x)=\dfrac{x^2+x-6}{|x-2|}$. (a) Tìm (i) $\lim_{x\to2^+}g(x)$, (ii) $\lim_{x\to2^-}g(x)$. (b) $\lim_{x\to2}g(x)$ có tồn tại không? (c) Vẽ phác đồ thị $g$.
\bt{52} Cho $f(x)=\begin{cases}x^2+1,&x<1,\\(x-2)^2,&x\ge1.\end{cases}$

(a) Tìm $\lim_{x\to1^-}f(x)$ và $\lim_{x\to1^+}f(x)$. (b) Giới hạn hai phía tại 1 có tồn tại không? (c) Vẽ phác đồ thị $f$.
\bt{53} Cho $B(t)=\begin{cases}4-\tfrac12t,&t<2,\\\sqrt{t+c},&t\ge2.\end{cases}$ Tìm $c$ để $\lim_{t\to2}B(t)$ tồn tại.
\bt{54} Cho $g(x)=\begin{cases}x,&x<1,\\3,&x=1,\\2-x^2,&1<x\le2,\\x-3,&x>2.\end{cases}$

(a) Tính nếu tồn tại:
\[(i)\ \lim_{x\to1^-}g(x);\quad(ii)\ \lim_{x\to1}g(x);\quad(iii)\ g(1);\]
\[(iv)\ \lim_{x\to2^-}g(x);\quad(v)\ \lim_{x\to2^+}g(x);\quad(vi)\ \lim_{x\to2}g(x).\]
(b) Vẽ phác đồ thị $g$.
\bt{55} (a) Với ký hiệu phần nguyên trong Ví dụ 10, tính
\[(i)\ \lim_{x\to-2^+}\gint{x};\quad(ii)\ \lim_{x\to-2}\gint{x};\quad(iii)\ \lim_{x\to-2.4}\gint{x}.\]
(b) Nếu $n$ nguyên, tính (i) $\lim_{x\to n^-}\gint{x}$, (ii) $\lim_{x\to n^+}\gint{x}$. (c) Với những $a$ nào, $\lim_{x\to a}\gint{x}$ tồn tại?
\bt{56} Cho $f(x)=\gint{\cos x}$ trên $-\pi\le x\le\pi$.

(a) Vẽ phác đồ thị. (b) Tính nếu tồn tại:
\[(i)\ \lim_{x\to0}f(x);\quad(ii)\ \lim_{x\to(\pi/2)^-}f(x);\]
\[(iii)\ \lim_{x\to(\pi/2)^+}f(x);\qquad(iv)\ \lim_{x\to\pi/2}f(x).\]
(c) Với những $a$ nào, $\lim_{x\to a}f(x)$ tồn tại?
\bt{57} Với $f(x)=\gint{x}+\gint{-x}$, chứng minh $\lim_{x\to2}f(x)$ tồn tại nhưng khác $f(2)$.
\bt{58} Trong thuyết tương đối, công thức co độ dài Lorentz
\[L=L_0\sqrt{1-v^2/c^2}\]
biểu diễn độ dài $L$ theo vận tốc $v$ đối với một người quan sát; $L_0$ là độ dài khi vật đứng yên, $c$ là tốc độ ánh sáng. Tìm $\lim_{v\to c^-}L$, giải thích kết quả và vì sao phải xét giới hạn bên trái.
\bt{59} Nếu $p$ là đa thức, chứng minh $\lim_{x\to a}p(x)=p(a)$.
\bt{60} Nếu $r$ là hàm hữu tỉ, dùng Bài tập 59 chứng minh $\lim_{x\to a}r(x)=r(a)$ với mọi $a$ thuộc tập xác định của $r$.
\bt{61} Nếu $\displaystyle\lim_{x\to1}\frac{f(x)-8}{x-1}=10$, tìm $\lim_{x\to1}f(x)$.
\bt{62} Nếu $\displaystyle\lim_{x\to0}\frac{f(x)}{x^2}=5$, tìm: (a) $\lim_{x\to0}f(x)$; (b) $\lim_{x\to0}f(x)/x$.
\bt{63} Cho $f(x)=\begin{cases}x^2,&x\text{ hữu tỉ},\\0,&x\text{ vô tỉ}.\end{cases}$ Chứng minh $\lim_{x\to0}f(x)=0$.
\bt{64} Nêu ví dụ cho thấy $\lim_{x\to a}[f(x)+g(x)]$ có thể tồn tại dù cả $\lim_{x\to a}f(x)$ và $\lim_{x\to a}g(x)$ đều không tồn tại.
\bt{65} Nêu ví dụ tương tự đối với $\lim_{x\to a}[f(x)g(x)]$.
\bt{66} Tính $\displaystyle\lim_{x\to2}\frac{\sqrt{6-x}-2}{\sqrt{3-x}-1}$.
\bt{67} Có số $a$ nào để $\displaystyle\lim_{x\to-2}\frac{3x^2+ax+a+3}{x^2+x-2}$ tồn tại không? Nếu có, tìm $a$ và giới hạn.
\bt{68} Hình vẽ cho đường tròn cố định $C_1:(x-1)^2+y^2=1$ và đường tròn $C_2$ tâm tại gốc tọa độ, bán kính $r$, đang co lại. Điểm $P=(0,r)$; $Q$ là giao điểm phía trên của hai đường tròn; $R$ là giao điểm đường thẳng $PQ$ với trục $x$. Điều gì xảy ra với $R$ khi $C_2$ co lại, tức $r\to0^+$?
\fig[.66]{e68}{Hình cho Bài tập 68.}
\section*{Ghi chú biên tập}
\begin{enumerate}
\item Ký hiệu ngoặc kép $\gint{x}$ chỉ hàm phần nguyên như bản gốc; tương đương $\lfloor x\rfloor$, không phải làm tròn về 0.
\item Quy tắc 7 và 11 giữ điều kiện dương nghiêm ngặt của sách khi bậc căn chẵn. Giới hạn căn tại đầu mút 0 trong Ví dụ 9 được xét một phía, không suy ra trực tiếp từ trường hợp hai phía với $a>0$.
\item Trong Ví dụ 11, các bất đẳng thức chứa $\sin(1/x)$ chỉ xét $x\ne0$. Chỉ cần các bất đẳng thức đúng gần điểm lấy giới hạn, không cần đúng tại chính điểm ấy.
\item Trong Bài tập 33, 34, $x$ là tham số cố định; riêng Bài tập 34 cần $x\ne0$. Các bài tập được dịch đề, không bổ sung đáp án.
\end{enumerate}
\section*{Thuật ngữ bổ sung}
\begin{center}\begin{tabular}{ll}\toprule Tiếng Anh&Tiếng Việt thống nhất\\\midrule
limit laws&các quy tắc giới hạn\\
sum law&quy tắc tổng\\
difference law&quy tắc hiệu\\
constant multiple law&quy tắc nhân với hằng số\\
product law&quy tắc tích\\
quotient law&quy tắc thương\\
power law&quy tắc lũy thừa\\
root law&quy tắc căn\\
direct substitution&phép thế trực tiếp\\
continuous at a point&liên tục tại một điểm\\
rationalizing the numerator&hữu tỉ hóa tử\\
squeeze theorem&định lý kẹp\\
greatest integer function&hàm phần nguyên\\
signum function&hàm dấu\\
Lorentz contraction&sự co độ dài Lorentz\\
method of exhaustion&phương pháp vét kiệt\\
\bottomrule\end{tabular}\end{center}
\noindent\textbf{Phạm vi hoàn thành:} Trọn §2.3, gồm 11 quy tắc giới hạn, tính chất thế trực tiếp, Định lý 1–3, Ví dụ 1–11, Hình 1–8, công thức (4) và 68 bài tập. Phần tiếp theo: \textit{§2.4 — Định nghĩa chính xác của giới hạn}, trang in 105.
\end{document}
